% ---------------------------------------------------------------------------
% Progress report section, appended below the original research plan.
% The text above this point (Introduction and proposed plan) is unchanged.
% All numbers come from the reference runs in the project repository.
% ---------------------------------------------------------------------------
\clearpage
\section{Progress report (November 2026)}

This section reports what has been implemented and found so far, following the plan above.
Every number comes from reproducible reference runs of the pipeline. \textbf{All results are
computational candidate relationships ranked by the available evidence. None demonstrates
biological equivalence, and none has been tested experimentally.}

\subsection{Implementation (Phase 1)}

\begin{itemize}
  \item \textbf{Networks.} \emph{D.~radiodurans} (NCBI 1299, mapped by STRING to strain R1,
    243230) and human (9606) networks are built from STRING v12.5 at combined score
    $\geq 0.7$ (configurable). Functional associations and STRING's physical subnetwork are
    kept apart, and experimentally confirmed physical interactions are checked separately
    in IntAct.
  \item \textbf{Seeds.} Proteins are selected from UniProtKB by GO annotation (including
    descendant terms): oxidative stress (GO:0006979, GO:0098869, GO:0016209) and DNA repair
    (GO:0006281). Seeds without a STRING entry are reported, never substituted.
  \item \textbf{Annotation.} GO with evidence codes (QuickGO; NOT qualifiers are treated as
    contradictions), KEGG pathways and KEGG Orthology, EC numbers, PANTHER, Pfam and InterPro.
  \item \textbf{Orthology as a feature, not a filter.} Instead of OrthoFinder, homology is graded
    in tiers: reciprocal best hit, shared KEGG KO, shared PANTHER family, one-way best hit,
    shared Pfam, shared InterPro, and a ``common fold'' tier for frequent fold-level signatures.
  \item \textbf{Reproducibility.} All API responses are cached immutably with SHA-256
    checksums, every run writes a manifest, the full analysis reruns offline, and result tables
    are byte-identical across reruns. The pipeline runs with one command and has 88 automated
    tests with mocked APIs.
\end{itemize}

\subsection{Oxidative-stress pilot}

\begin{table}[htbp]
\centering
\small
\begin{tabular}{lrr}
\toprule
 & \emph{D. radiodurans} & \emph{H. sapiens} \\
\midrule
proteins in network & 16 & 382 \\
edges (score $\geq 0.7$) & 11 & 2,210 \\
isolated proteins & 9 & -- \\
IntAct-confirmed STRING edges & 0 of 11 & 275 of 2,210 (12.4\%) \\
\bottomrule
\end{tabular}
\caption{Oxidative-stress networks.}
\end{table}

\begin{itemize}
  \item \textbf{Coverage.} The bacterial network is small because few proteins carry the seed
    GO terms, not because they are poorly connected: random 16-protein human subsets have on
    average 3.5 edges, the bacterial set has 11. Adding every bacterial gene sharing a KEGG KO
    with a human seed only grows the network to 27 proteins and 16 edges.
  \item \textbf{Modules.} Louvain (stability over 50 seeds, mean ARI 0.74 for human) finds a
    69-protein human antioxidant module (SOD1--3, CAT, PRDX1--6, GPX2--8, TXNRD1--3,
    NFE2L2/KEAP1), plus proteasome and mitochondrial modules. Both bacterial modules match the
    human antioxidant module best.
  \item \textbf{Alignment.} Global alignment (IsoRank) conserves more edges than random
    mappings, but not more than degree-preserving rewired networks (p = 0.10 without prior;
    0.63--0.88 with a homology prior). \textbf{Topology alone shows no conserved wiring.}
  \item \textbf{Orthologs (sanity check).} Reciprocal best hits recover katA--CAT, sodA--SOD2,
    msrB--MSRB3 and the bacterial peroxiredoxin with PRDX2.
  \item \textbf{Candidate analogues.} Superoxide dismutation by unrelated families: bacterial
    Mn/Fe-SOD (sodA) with human Cu/Zn-SOD3 and SOD1, and bacterial Cu/Zn proteins with human
    Mn-SOD2. Peroxide removal by catalase (katA, DRA0259) versus peroxiredoxins is an analogy of
    reaction class only.
  \item \textbf{Mixed evidence.} sodA mutants are more radiation-sensitive in one study
    \cite{markillie1999}, but $\Delta$MnSOD mutants survive gamma radiation better in another,
    which attributes protection to small-molecule Mn$^{2+}$ antioxidants
    \cite{gaidamakova2022}. A protein network cannot represent these non-protein defences.
  \item \textbf{Worth an assay.} OsmC (DR1538) pairs with human GPX2/7/8 on shared function
    and a curated disruption phenotype \cite{wang2016}, without any detectable homology.
\end{itemize}

\subsection{DNA repair}

\begin{table}[htbp]
\centering
\small
\begin{tabular}{lrr}
\toprule
 & \emph{D. radiodurans} & \emph{H. sapiens} \\
\midrule
seeds mapped to STRING & 52 of 56 & 426 of 432 \\
edges (score $\geq 0.7$) & 119 & 5,285 \\
largest component & 41 & 400 \\
IntAct-confirmed STRING edges & 0 of 119 & 431 of 5,285 (8.2\%) \\
\bottomrule
\end{tabular}
\caption{DNA-repair networks.}
\end{table}

\begin{itemize}
  \item \textbf{Bacterial modules recover known pathways.} All five non-singleton Louvain
    modules are FDR-significant: base-excision repair, RecFOR recombination,
    nucleotide-excision repair, replication/mismatch repair, and a radiation-response module
    (recA, lexA, ddrA, ddrB, pprA).
  \item \textbf{Module correspondence is weak.} Several bacterial modules match the human
    base-excision module best, because generic repair terms dominate the module profiles.
  \item \textbf{Alignment.} Without a prior there is no significant conservation (p = 0.08 against
    rewired networks). Anchored on homologs, conservation exceeds chance (p = 0.01--0.02), but STRING
    transfers evidence between orthologs, so this is not independent evidence of conserved
    wiring.
  \item \textbf{Orthologs.} nth--NTHL1, DR\_0354--APEX1, DR\_2285--MUTYH, ung--UNG,
    mutL--MLH1, recA--DMC1; homologs include recA--RAD51, RecQ--WRN and ddrA--RAD52
    \cite{harris2004}.
  \item \textbf{Convergent glycosylases (clearest analogues).} mutM (Fpg family) with NTHL1 and
    OGG1 (HhH-GPD family), and nth with NEIL1/NEIL3: different folds, same AP-lyase activity.
    Bacterial MutM and human OGG1 both remove 8-oxoguanine.
  \item \textbf{Pathway-level analogy.} UvrABC versus XPC/TFIIH in nucleotide-excision repair,
    and recO versus RAD52 in recombination: same pathway, largely different proteins.
  \item \textbf{Curated evidence changes interpretations.} lexA deletion does not change
    radiation survival \cite{jolivet2006}, despite its GO ``DNA repair'' label; pprA loss
    reduces radiation resistance \cite{narumi2004}, but $\Delta$pprA repairs double-strand
    breaks efficiently and its defect lies in segregation \cite{devigne2013}. pprA has no
    supported human counterpart.
\end{itemize}

\subsection{Ranking and robustness (Phase 4)}

Each pair gets a transparent weighted mean over the evidence components available for it
(EC, GO, KEGG pathway, literature, perturbation, module, network role, homology), so missing
evidence never counts as negative. Pairs with too little evidence are listed but not ranked.
Pairs are labelled ortholog, homolog, distant homolog, candidate analogue, shared function,
similar pathway role or speculative.

\begin{table}[H]
\centering
\small
\begin{tabular}{L{4.2cm}L{3.2cm}L{3.2cm}}
\toprule
check & oxidative stress & DNA repair \\
\midrule
500 random weight sets: top-20 overlap (median) & 0.80 & 0.70 \\
STRING threshold 0.4--0.9: top-20 overlap & 0.80--1.00 & 0.60--1.00 \\
candidate analogues (3 GO similarity measures) & 65 / 81 / 120 & 67 / 90 / 108 \\
NetMF link-prediction AUC (human) & 0.935 (Adamic--Adar 0.949) & 0.879 (Adamic--Adar 0.952) \\
\bottomrule
\end{tabular}
\caption{Sensitivity checks. Embeddings did not beat simple neighbourhood baselines and are
not used in the ranking.}
\end{table}

There is no ground truth for cross-species analogy, so no accuracy is claimed. Network-role
similarity alone does not recover known orthologs (AUC 0.46 on 16 DNA-repair ortholog pairs).

\subsection{Status against the proposed plan}

\begin{table}[H]
\centering
\small
\begin{tabular}{L{0.9cm}L{4.9cm}L{4.9cm}}
\toprule
phase & done & not done / different \\
\midrule
1 & STRING $\geq 0.7$ networks; oxidative stress and DNA repair; KEGG/GO; orthology as a feature
  & broader stress-response subsystem; OrthoFinder (replaced by RBH, KEGG KO, PANTHER, Pfam/InterPro) \\
2 & degree, betweenness, closeness, clustering, hub/bottleneck; Louvain with stability
  & graphlet degree vectors; Leiden; transcriptomic overlay of irradiation time courses \\
3 & module correspondence by enrichment profiles; global alignment seeded on homology, with
  rewired nulls
  & module-preservation statistics; MAGNA++ or graphlet alignment; cross-species embeddings \\
4 & transparent weighted score, evidence categories, sensitivity analysis & -- \\
5 & candidate shortlist with proposed experiments; curated bacterial mutant evidence
  & cross-check against human genetics (constraint, disease associations) \\
\bottomrule
\end{tabular}
\caption{Progress against the five phases.}
\end{table}

\subsection{Limitations}

\begin{itemize}
  \item STRING edges are functional associations, partly transferred between organisms through
    orthology. IntAct, BioGRID and MINT hold no physical interactions for the bacterial seed
    proteins, so physical-interaction comparison across the two species is not possible.
  \item Most bacterial GO annotation is electronic. ``No detected homology'' cannot exclude
    very remote homology.
  \item Curated mutant evidence covers 12 bacterial proteins, from abstracts only. Human
    perturbation evidence comes only from GO evidence codes.
  \item Broad EC numbers (e.g. EC 2.7.7.7, any DNA polymerase) inflate some analogues.
\end{itemize}

\subsection{Proposed next steps}

\begin{enumerate}
  \item Cross-complementation tests: \emph{D.~radiodurans} sodA deletion with human SOD1/SOD3
    (under H$_2$O$_2$ and desiccation), mutM deletion with human OGG1, and nth with NEIL1.
  \item An activity assay for OsmC (DR1538) to settle the comparison with GPXs.
  \item The missing Phase 2--3 methods: graphlet degree vectors, Leiden, module-preservation
    statistics and a local aligner, then cross-species embeddings.
  \item Cross-check the shortlist against human constraint and disease-association data.
  \item Represent Mn$^{2+}$--metabolite antioxidants explicitly as network context.
\end{enumerate}

\subsection{Reproducibility}

Code, cached data, reference runs and detailed reports are in the project repository.

\begin{itemize}
  \item Repository: \url{https://github.com/Ananthapadmanabhan333/Cross-Species-Oxidative-States-}
    (release \texttt{v0.3.0})
  \item Oxidative-stress reference run: \texttt{20261103T111919Z\_860ea5376e14}
  \item DNA-repair reference run: \texttt{20261103T112433Z\_73788e2787df}
\end{itemize}

\begin{thebibliography}{9}
\bibitem{markillie1999} Markillie et al. (1999), \href{https://doi.org/10.1128/JB.181.2.666-669.1999}{doi:10.1128/JB.181.2.666-669.1999}.
\bibitem{gaidamakova2022} Gaidamakova et al. (2022), \href{https://doi.org/10.1128/mbio.03394-21}{doi:10.1128/mbio.03394-21}.
\bibitem{wang2016} Wang et al. (2016), \href{https://doi.org/10.3109/09553002.2016.1150618}{doi:10.3109/09553002.2016.1150618}.
\bibitem{harris2004} Harris et al. (2004), \href{https://doi.org/10.1371/journal.pbio.0020304}{doi:10.1371/journal.pbio.0020304}.
\bibitem{jolivet2006} Jolivet et al. (2006), \href{https://doi.org/10.1111/j.1365-2958.2005.04946.x}{doi:10.1111/j.1365-2958.2005.04946.x}.
\bibitem{narumi2004} Narumi et al. (2004), \href{https://doi.org/10.1111/j.1365-2958.2004.04272.x}{doi:10.1111/j.1365-2958.2004.04272.x}.
\bibitem{devigne2013} Devigne et al. (2013), \href{https://doi.org/10.1016/j.dnarep.2013.01.004}{doi:10.1016/j.dnarep.2013.01.004}.
\end{thebibliography}
